\section{Introduzione}

\begin{paragrafo}{Scopo}
    L'obiettivo della materia è quello di andare a studiare la complessità degli algoritmi, andare ad approssimare algoritmi che non sono
    applicabili, e studiare queste approssimazioni.

    Aggiungo inoltre che dato che ho appena iniziato a farla, non so niente, e questo paragrafo verrà cambiato nel futuro lel. %TODO cambiarlo.
\end{paragrafo}

\subsection{Notazione e cenni matematici}

    In questa sezione ci saranno un po' di cenni ai concetti che andremo a utilizzare in tutto il corso e le notazioni usate.
    \separapassaggi
    \begin{paragrafo}{Alfabeto}
        Un insieme finito e non vuoto è detto alfabeto e si indica generalmente con $\Sigma$
    \end{paragrafo}

    \begin{paragrafo}{Insieme delle stringhe}
        L'insieme delle stringhe di un alfabeto è l'insieme di tutte le possibili sequenze finite su quell'alfabeto e si indica con $\Sigma^*$
    \end{paragrafo}
    
    \begin{paragrafo}{Numerabilità}
        Un insieme è detto \textbf{numberabile} se questo  può essere messo in biiezione con i naturali.
    \end{paragrafo}

    \begin{paragrafo}{Concatenazione}
        L'operazione di concatenazione, è un operatore che date due stringhe $s1,s2 \in \Sigma^*$ crea la stringa $s1 \circ s2$, dove $\circ$ concatena le stringhe.

        L'operatore di concatenazione è associativa e non commutativa.
    \end{paragrafo}

    \begin{paragrafo}{Magma, Semigruppo, Monoide }
        Un magma è una coppia insieme-operazione, in cui l'operazione non è associativa

        Un semigruppo è una coppia iniseme-operazione, in cui l'operazione è associativa $(\Sigma^*, \circ)$ 

        Un monoide è un semigruppo in cui l'operazione ha elemento neutro nell'insieme.

        Un monoide abeliano è un monoide in cui l'operazione gode anche della proprietà commutativa.
    \end{paragrafo}
    
    \begin{paragrafo}{Insime delle funzioni}
        L'insieme delle funzioni $B^A$ è l'insieme di tutte le possibili funzioni con dominio in $A$ e codominio in $B$

        \[
            B^A = \{f\}\ f:A\rightarrow B
        \]

    \end{paragrafo}

    \begin{paragrafo}{Naturali come insiemi}
        Spesso nel corso capiterà di utilizzare i numeri naturali per andare a indicare il sottoinsieme dei naturali fino a quel numero escluso.
        
    \[
        0 = \emptyset \qquad 1 = \{0\} \qquad 2 = \{0,1\} \qquad 3 = \{0,1,2\}
    \]

    \[
        2^* = \{\epsilon,0,1,00,01,11,10,\dots \} \qquad 2^A= \{f\}\ f:A \rightarrow 2
    \]

    \[
        2^{2^*} = \text{ Famiglia che descrive tutti i linguaggi binari}
    \]
    \begin{center}che non è numerabile, per il teorema di Cantor\end{center}

    \[
        A^2=\{(a_1,a_2) : a_1, a_2 \in A\} = \{f:2 \rightarrow A\} = \{f:\{0,1\}\rightarrow A\}
    \]

    \end{paragrafo}


    \begin{paragrafo}{Problemi di decisione} 
        Dato $L \subseteq \Sigma^*$, un problema di decisione è un problema che dato in input un elemento di $\Sigma^*$ da in output si o no, a seconda dell'appartenenza dell'elemento a $L$

    \[
        x\in \Sigma^* \rightarrow \boxed{P_L} \rightarrow
        \begin{cases}
            1 & x \in L \\
            0 & x \notin L 
        \end{cases}
    \]

    \end{paragrafo}
    
    \begin{paragrafo}{Problema}
        Un problema $\pi$ è definito da:
            \begin{itemize}
                \item un insieme di input possibili $I_\pi \subseteq 2^*$
                \item un insieme di output possibili $O_\pi \subseteq 2^*$
                \item una funzione $Sol_\pi: I_\pi\rightarrow 2^{O_\pi}$
            \end{itemize}
    \end{paragrafo}

    \begin{paragrafo}{Rappresentazione degli input}
        Gli input di un problema come descritti sopra devono essere un unica stringa binaria. Data questa limitazione, è necessario trovare un modo per andare a codificare più numeri di input in una singola stringa in input.
        Questo viene fatto tramite alla \textbf{gamma di Elios}, che consiste nel codificare i numeri in binario andando a inserire ad ognuno un prefisso composto tanti 1 quanto le cifre del numero, e uno 0 per indicare la fine del prefisso.
        Questo sicuramente aumenta la lunghezza dell'input, ma lo fa di un valore irrisorio all'ordine di grandezza, quindi trascurabile.
    \end{paragrafo}


    \begin{paragrafo}{Soluzioni}
        Le soluzioni sono dettate dagli algoritmi e sono così definite:
        \[
            x\in I_\pi \rightarrow \boxed{A}_\pi \rightarrow y \in O_\pi \land y \in Sol_\pi(x)
        \]
    \end{paragrafo}

    \begin{paragrafo}{Algoritmo}
        Un algoritmo è un insieme di istruzioni, noi li vedremo come macchine di Turing.
    \end{paragrafo}

    \begin{paragrafo}{Macchina di Turing}
        Una macchina di Turing è una macchina ideale, composta da un nastro infinito con su scritti dati e separatori.

        Un puntatore scorre sul nastro avanti e indietro e calcola dai dati scritti una soluzione.

        \begin{paragrafo}{Tesi di Church-Turing}
            La tesi di Church-Turing dice che la macchina di Turing rappresenta il massimo della potenza computazionale ad oggi raggiungibile. Ogni altro strumento di calcolo, come i moderni linguaggi, è pari ad essa.
        \end{paragrafo}
    \end{paragrafo}

    % TODO finire 
    

\paragraph{}
